\chapter{曲线：从边界条件到数据拟合}
\label{ch:curves}

\section{参数表达与三次Hermite曲线}
\sourcepage{第4章}{3--8}

隐式表达适合判定点是否位于曲线上，例如圆的 $x^2+y^2-R^2=0$；显式表达 $y=f(x)$ 则不能用一条单值函数完整描述圆，也不便处理竖直切线。CAD更常用参数表达
\[
 \bs C(u)=(x(u),y(u),z(u)),\qquad u\in[a,b].
\]
参数 $u$ 用来定位曲线上的点，通常没有长度单位，也不必等于时间或弧长。参数范围给出曲线段的边界；对 $\bs C$ 的各坐标同时作仿射变换，就得到变换后的曲线。几何形状相同的曲线可以有不同的参数化，因此比较导数时必须先统一参数。

\dfn{三次Hermite曲线}{给定端点 $\bs P_0,\bs P_1$ 及关于 $u$ 的端点导矢 $\bs T_0,\bs T_1$，求满足
\[
 \bs C(0)=\bs P_0,\quad \bs C(1)=\bs P_1,\quad
 \bs C'(0)=\bs T_0,\quad\bs C'(1)=\bs T_1
\]
的三次参数曲线。四个边界条件分别决定三个坐标分量的四个系数。}

令 $\bs C(u)=\bs a_0+\bs a_1u+\bs a_2u^2+\bs a_3u^3$。代入边界条件，先得 $\bs a_0=\bs P_0,\bs a_1=\bs T_0$，再解
\[
 \bs a_2+\bs a_3=\bs P_1-\bs P_0-\bs T_0,\quad
 2\bs a_2+3\bs a_3=\bs T_1-\bs T_0,
\]
于是
\begin{align*}
 \bs a_2&=-3\bs P_0+3\bs P_1-2\bs T_0-\bs T_1,\\
 \bs a_3&=2\bs P_0-2\bs P_1+\bs T_0+\bs T_1.
\end{align*}
整理为几何形式：
\begin{align*}
 \bs C(u)&=h_{00}(u)\bs P_0+h_{01}(u)\bs P_1
             +h_{10}(u)\bs T_0+h_{11}(u)\bs T_1,\\
 h_{00}&=1-3u^2+2u^3,& h_{01}&=3u^2-2u^3,\\
 h_{10}&=u-2u^2+u^3,& h_{11}&=-u^2+u^3.
\end{align*}
这些混合函数将边界条件与曲线形状联系起来：端点不变时，改变导矢方向会改变出发与到达方向，改变导矢模长则改变端部的牵引程度。

\reviewfig[.64]{ch04-hermite-tangents.pdf}{端点相同而端部导矢不同的Hermite曲线。第4章讲义，PDF第8页。}{fig:hermite}

\ex{用边界条件检查结果}{取 $\bs P_0=(0,0)$、$\bs P_1=(2,0)$、$\bs T_0=(2,2)$、$\bs T_1=(2,-2)$，代入得到
\[
 \bs C(u)=(2u,\,2u-2u^2),\qquad \bs C(1/2)=(1,1/2).
\]
这里三次项恰好消失，实际次数降为2。检验端点和导数比只核对展开系数更直接。}

\nt{三次多项式具有独立指定两个端点及两个导矢的能力，但任意两段三次曲线不会自动达到 $C^2$ 连续；连续性仍需施加拼接条件。讲义中Hermite系数应以四个边界条件检验，上式保留了末端导矢的正确符号。}

\section{Bézier曲线：控制多边形与几何运算}
\sourcepage{第4章}{9--22}

\subsection{定义与基本性质}
\dfn{Bernstein基函数与Bézier曲线}{由 $n+1$ 个控制点 $\bs P_0,\ldots,\bs P_n$ 定义
\[
 \bs C(u)=\sum_{i=0}^{n}B_i^n(u)\bs P_i,\quad
 B_i^n(u)=\binom ni u^i(1-u)^{n-i},\quad 0\leq u\leq1.
\]
$\bs P_i$ 顺次相连形成控制多边形。曲线最高为 $n$ 次；特定控制点配置可能使实际次数降低。}

由二项式定理，$B_i^n(u)\geq0$ 且 $\sum_iB_i^n(u)=1$。因此每个曲线点都是控制点的凸组合，直接得到以下性质。
\begin{itemize}
 \item \textbf{凸包性：}整条曲线位于控制点的凸包内，可用于包围盒测试和求交剔除。
 \item \textbf{端点插值：}$\bs C(0)=\bs P_0,\bs C(1)=\bs P_n$；一般不经过中间控制点。
 \item \textbf{仿射不变性：}先变换控制点再求曲线，等于先求曲线再作同一仿射变换。平移项成立的关键正是单位分割。
 \item \textbf{全局修改：}在 $0<u<1$ 内各基函数均为正，移动一个控制点通常影响整段曲线。增加控制点会提高整段曲线的次数，因此复杂形状常采用多段低次曲线。
\end{itemize}

对基函数求导并按控制点差分重组：
\begin{align*}
 \bs C'(u)&=n\sum_{i=0}^{n-1}B_i^{n-1}(u)(\bs P_{i+1}-\bs P_i),\\
 \bs C'(0)&=n(\bs P_1-\bs P_0),\quad
 \bs C'(1)=n(\bs P_n-\bs P_{n-1}),\\
 \bs C''(u)&=n(n-1)\sum_{i=0}^{n-2}B_i^{n-2}(u)
 (\bs P_{i+2}-2\bs P_{i+1}+\bs P_i).
\end{align*}
导曲线仍为Bézier形式，其控制点由原控制点的一阶或二阶差分得到。端点导数的\textbf{方向和模长}都由相邻控制点决定，不能仅画一条切线就认为导数已确定。

三次Bézier曲线与Hermite曲线描述的是同一类三次多项式。若Hermite端点导矢为 $\bs T_0,\bs T_1$，则等价控制点是
\[
 \bs B_0=\bs P_0,\quad \bs B_1=\bs P_0+\frac{\bs T_0}{3},\quad
 \bs B_2=\bs P_1-\frac{\bs T_1}{3},\quad\bs B_3=\bs P_1.
\]

\subsection{拼接：形状光滑与参数光滑}
设前段控制点为 $\bs P_0,\ldots,\bs P_n$，后段为 $\bs Q_0,\ldots,\bs Q_m$，各自参数区间均为 $[0,1]$。
\begin{align*}
 C^0\;(G^0):&\quad \bs P_n=\bs Q_0;\\
 C^1:&\quad n(\bs P_n-\bs P_{n-1})=m(\bs Q_1-\bs Q_0);\\
 G^1:&\quad n(\bs P_n-\bs P_{n-1})
       =\lambda m(\bs Q_1-\bs Q_0),\quad\lambda>0.
\end{align*}
$C^1$ 要求导矢相等；$G^1$ 在端点导矢非零时只要求同向平行，允许参数速度不同。两段次数相同且参数尺度相同，$C^1$ 才简化成接点两边的控制边矢量相等。

若还要求 $C^2$，则须进一步满足
\[
 n(n-1)(\bs P_n-2\bs P_{n-1}+\bs P_{n-2})
 =m(m-1)(\bs Q_2-2\bs Q_1+\bs Q_0).
\]
若两段对应的全局参数跨度为 $\Delta u_1,\Delta u_2$，一阶导数还应分别除以跨度，二阶导数分别除以跨度的平方。

\ex{三次曲线的接点控制}{前段末端 $\bs P_2=(1,1),\bs P_3=(2,0)$，后段起点取 $\bs Q_0=(2,0)$。要达到 $C^1$，应取
\[
 \bs Q_1=\bs Q_0+(\bs P_3-\bs P_2)=(3,-1).
\]
若改成 $\bs Q_1=(4,-2)$，导矢同向但模长变为两倍，只保证 $G^1$。若取 $\bs Q_1=(1,1)$，导矢反向，不能按同向行进的 $G^1$ 光滑拼接处理。}

\subsection{升阶：表示改变而曲线不变}
由 $\bs C=u\bs C+(1-u)\bs C$，将两部分均整理成 $n+1$ 次Bernstein基，可得新控制点
\begin{align*}
 \bs Q_0&=\bs P_0,\qquad \bs Q_{n+1}=\bs P_n,\\
 \bs Q_i&=\frac{i}{n+1}\bs P_{i-1}
          +\left(1-\frac{i}{n+1}\right)\bs P_i,\quad 1\leq i\leq n.
\end{align*}
升阶增加控制点和表达次数，\textbf{曲线轨迹及参数化均不变}。它适合统一两段曲线的次数，不等于增添新的形状自由度后任意移动控制点。

\ex{二次升为三次}{原控制点为 $(0,0),(1,2),(2,0)$，升阶后为
\[
 (0,0),\quad(2/3,4/3),\quad(4/3,4/3),\quad(2,0).
\]
用两种表达计算 $u=1/2$，均得 $(1,1)$。}

\subsection{de Casteljau求值与细分}
直接展开高次多项式既不便计算，也不直观。de Casteljau算法在相邻点间反复作同一比例的线性插值：
\[
 \bs P_i^{(0)}=\bs P_i,\qquad
 \bs P_i^{(r)}=(1-u)\bs P_i^{(r-1)}+u\bs P_{i+1}^{(r-1)},
\]
其中 $r=1,\ldots,n$、$i=0,\ldots,n-r$，最后得到 $\bs C(u)=\bs P_0^{(n)}$。

固定细分参数 $u=u_*$ 后，递推三角形的左边界给出前半段控制点
\[
 \bs P_0^{(0)},\bs P_0^{(1)},\ldots,\bs P_0^{(n)},
\]
右边界反向给出后半段控制点
\[
 \bs P_0^{(n)},\bs P_1^{(n-1)},\ldots,\bs P_n^{(0)}.
\]
两段分别对应原参数区间 $[0,u_*]$ 与 $[u_*,1]$，各自重新归一化到 $[0,1]$ 后形状不变。

\reviewfig[.74]{ch04-casteljau-split.pdf}{de Casteljau递推同时完成求值与细分；两条边界上的点成为子曲线控制点。第4章讲义，PDF第22页。}{fig:casteljau}

\ex{三次曲线的递推计算}{取 $(0,0),(0,2),(2,2),(2,0)$，求 $u=1/2$：第一层为 $(0,1),(1,2),(2,1)$，第二层为 $(1/2,3/2),(3/2,3/2)$，第三层为 $(1,3/2)$。因此曲线点为 $(1,3/2)$，前半段控制点为
\[
 (0,0),\ (0,1),\ (1/2,3/2),\ (1,3/2).
\]
}

\section{B样条：用节点组织局部曲线}
\sourcepage{第4章}{23--44}

\subsection{次数、节点与Cox--de Boor递推}
Bézier将控制点数量与次数绑定；B样条引入\textbf{节点向量}，使多个低次多项式段共同表达一条曲线。以下统一用 $p$ 表示次数，$k=p+1$ 表示讲义的阶次，避免将三次曲线误写成 $k=3$。

\dfn{$p$次B样条曲线}{给定 $n+1$ 个控制点及非递减节点向量
\[
 U=(t_0,t_1,\ldots,t_{n+p+1}),\quad t_i\leq t_{i+1},
\]
曲线为 $\bs C(u)=\sum_{i=0}^nN_{i,p}(u)\bs P_i$，常用有效参数域为 $[t_p,t_{n+1}]$。这里 $N_{i,p}$ 按\textbf{次数}编号，对应讲义的 $N_{i,k}$，两者下标相差1。}

基函数递推为
\begin{align*}
 N_{i,0}(u)&=\begin{cases}1,&t_i\leq u<t_{i+1},\\0,&\text{其他},\end{cases}\\
 N_{i,p}(u)&=\frac{u-t_i}{t_{i+p}-t_i}N_{i,p-1}(u)
 +\frac{t_{i+p+1}-u}{t_{i+p+1}-t_{i+1}}N_{i+1,p-1}(u).
\end{align*}
因重节点而出现的零分母项整体按0处理。右端点按左极限定义；端点夹持时令 $N_{n,p}(t_{n+1})=1$，以保留末端插值。零次基函数使用半开区间，不能让同一非重节点同时落入两个区间。

\begin{itemize}
 \item \textbf{局部支撑：}$N_{i,p}$ 仅可能在 $[t_i,t_{i+p+1})$ 内非零。移动 $\bs P_i$ 只影响该参数范围。
 \item \textbf{局部凸包：}在非零节点区间 $[t_\ell,t_{\ell+1})$ 内，至多有 $p+1$ 个基函数非零，对应 $\bs P_{\ell-p},\ldots,\bs P_\ell$。
 \item \textbf{非负与单位分割：}在有效参数域内仍有 $N_{i,p}\geq0,\sum_iN_{i,p}=1$，所以保持凸包性与仿射不变性。
 \item \textbf{连续性：}次数为 $p$，内部节点重数为 $r\leq p$ 时，一般保证 $C^{p-r}$ 连续。单节点为 $C^{p-1}$；提高重数会降低连续性，特殊控制点配置可以更光滑。
\end{itemize}

\nt{节点的\textbf{整体平移或统一缩放}可由参数重标定吸收；内部节点之间的相对间隔会影响基函数和曲线形状，不能理解为“节点间距完全不重要”。节点重数也直接影响光滑程度。}

\subsection{夹持节点与Bézier特例}
端点各重复 $p+1$ 次称为端点夹持，例如二次曲线
\[
 U=(0,0,0,1,2,3,4,4,4),\quad p=2,\quad n=5.
\]
它通过 $\bs P_0$ 和 $\bs P_5$，包含4个非零节点区间；内部节点均为单节点，通常为 $C^1$ 连续。图\ref{fig:bspline-basis}显示每个控制点的作用范围。

\reviewfig[.72]{ch04-bspline-basis.pdf}{二次B样条的局部基函数。图中 $N_{i,3}$ 使用讲义的“阶次3”记号，对应正文的 $N_{i,2}$。第4章讲义，PDF第35页。}{fig:bspline-basis}

若仅有 $p+1$ 个控制点，节点为
\[
 U=(\underbrace{0,\ldots,0}_{p+1},\underbrace{1,\ldots,1}_{p+1}),
\]
则基函数就是Bernstein基，B样条退化为单段Bézier曲线。\textbf{均匀}描述非零节点间隔是否相同，\textbf{夹持}描述端点重数，两种性质应分别判断。未夹持的均匀节点通常不插值两端控制点；若要构造闭合周期曲线，还需控制点和首尾导数满足周期条件，不能仅凭“均匀节点”断言曲线闭合。

\subsection{de Boor求值}
求某一 $u\in[t_\ell,t_{\ell+1})$ 的曲线点时，只使用局部 $p+1$ 个控制点，令
\[
 \bs D_i^{(0)}=\bs P_i\quad(i=\ell-p,\ldots,\ell).
\]
逐层递推
\begin{align*}
 \alpha_{i,r}&=\frac{u-t_i}{t_{i+p-r+1}-t_i},\\
 \bs D_i^{(r)}&=(1-\alpha_{i,r})\bs D_{i-1}^{(r-1)}
                 +\alpha_{i,r}\bs D_i^{(r-1)},\\
 r&=1,\ldots,p,\qquad i=\ell-p+r,\ldots,\ell.
\end{align*}
最终 $\bs C(u)=\bs D_\ell^{(p)}$。与de Casteljau相比，结构同为逐层插值，但插值比例由\textbf{当前节点间距}决定，通常不同。

\ex{讲义节点向量上的 $u=2.5$}{取 $p=2$ 和上面的 $U$，$u=2.5\in[t_4,t_5)$，所以 $\ell=4$，仅用 $\bs P_2,\bs P_3,\bs P_4$。
\begin{align*}
 \bs D_3^{(1)}&=\tfrac14\bs P_2+\tfrac34\bs P_3,\\
 \bs D_4^{(1)}&=\tfrac34\bs P_3+\tfrac14\bs P_4,\\
 \bs C(2.5)&=\tfrac12\bs D_3^{(1)}+\tfrac12\bs D_4^{(1)}
             =\tfrac18\bs P_2+\tfrac34\bs P_3+\tfrac18\bs P_4.
\end{align*}
系数非负且和为1，可作为快速检验。若再取 $\bs P_2=(0,0),\bs P_3=(2,2),\bs P_4=(4,0)$，则得到 $(2,1.5)$。}

\subsection{插入节点与分段编辑}
在 $\xi\in[t_\ell,t_{\ell+1})$ 插入一个节点，保持次数和曲线不变，同时增加一个控制点。若 $\xi$ 原重数为 $s$，一次插入后的控制点满足
\[
 \bs Q_i=\begin{cases}
 \bs P_i,&0\leq i\leq\ell-p,\\
 (1-\alpha_i)\bs P_{i-1}+\alpha_i\bs P_i,
     &\ell-p+1\leq i\leq\ell-s,\\
 \bs P_{i-1},&\ell-s+1\leq i\leq n+1,
 \end{cases}
 \quad \alpha_i=\frac{\xi-t_i}{t_{i+p}-t_i}.
\]
插入过程重新分配控制点而不改变几何形状；随后移动新增控制点才会改变曲线。反复插入内部节点至重数 $p$ 可将曲线表达为相接Bézier段。\textbf{删除节点}反向减少表示复杂度，但一般只能在给定误差容限内近似删除，不能保证任意节点都可无损移除。

\section{NURBS：权重与圆锥曲线}
\sourcepage{第4章}{45--49}

\dfn{非均匀有理B样条}{对B样条控制点增加权重 $w_i$，得到
\[
 \bs C(u)=\frac{\sum_{i=0}^nw_iN_{i,p}(u)\bs P_i}
                 {\sum_{i=0}^nw_iN_{i,p}(u)}.
\]
通常取 $w_i>0$，保证分母非零并保持凸包性。先在四维齐次空间对 $(w_i\bs P_i,w_i)$ 作普通B样条计算，再除以最后一维，就得到三维曲线。}

所有权重相等时退化为普通B样条；将\textbf{全部权重乘以同一非零常数}不改变曲线。只增加某个正权重时，在其基函数非零处，曲线点向相应控制点靠近；作用范围仍由B样条基函数的局部支撑决定。若端点夹持且端点权重非零，曲线仍通过首末控制点。

普通多项式Bézier或B样条不能精确表达非退化圆弧，\textbf{有理二次}曲线可以。对小于 $180^\circ$ 的圆弧，令 $\bs P_0,\bs P_2$ 为端点、$\bs P_1$ 为两端切线交点，圆心角为 $2\theta$，取
\[
 p=2,\quad U=(0,0,0,1,1,1),\quad
 (w_0,w_1,w_2)=(1,\cos\theta,1).
\]
$180^\circ$ 时两端切线平行，这种\textbf{有限控制点}构造退化，宜拆成两段或更多段；大圆弧也可分段拼合。

\reviewfig[.40]{ch04-nurbs-arc.pdf}{圆弧的两端点与切线交点决定有理二次控制三角形。第4章讲义，PDF第47页，裁取几何示意。}{fig:nurbs-arc}

\ex{单位圆四分之一弧}{取 $\bs P_0=(1,0)$、$\bs P_1=(1,1)$、$\bs P_2=(0,1)$，权重为 $(1,1/\sqrt2,1)$。在 $u=1/2$，二次Bernstein系数为 $(1/4,1/2,1/4)$，因此
\[
 \bs C(1/2)=
 \frac{(\tfrac14+\tfrac{1}{2\sqrt2},\,
              \tfrac14+\tfrac{1}{2\sqrt2})}
      {\tfrac12+\tfrac{1}{2\sqrt2}}
 =(1/\sqrt2,1/\sqrt2).
\]
若把中间权重改为1，则得到 $(3/4,3/4)$，不在单位圆上。}

\section{曲线求交：全局定位与局部迭代}
\sourcepage{第4章}{50--56}

对两条平面参数曲线 $\bs P(u),\bs Q(v)$，交点满足
\[
 \bs F(u,v)=\begin{pmatrix}P_x(u)-Q_x(v)\\P_y(u)-Q_y(v)\end{pmatrix}=\bs0.
\]
在当前参数 $(u_j,v_j)$ 处作一阶展开，求解
\[
 \underbrace{\begin{pmatrix}P_x'(u_j)&-Q_x'(v_j)\\P_y'(u_j)&-Q_y'(v_j)\end{pmatrix}}_{J}
 \begin{pmatrix}\Delta u\\\Delta v\end{pmatrix}
 =-\bs F(u_j,v_j),\qquad
 (u_{j+1},v_{j+1})=(u_j+\Delta u,v_j+\Delta v).
\]
这就是二维Newton迭代。应同时检查参数范围、参数增量与几何残差 $\|\bs P(u)-\bs Q(v)\|$。空间曲线还要验证第三个坐标，平面投影相交不能保证空间相交。

\ex{一次线性迭代}{设 $\bs P(u)=(u,u)$、$\bs Q(v)=(v,1-v)$，则 $\bs F=(u-v,u+v-1)$、$J=\left(\begin{smallmatrix}1&-1\\1&1\end{smallmatrix}\right)$。从 $(0,0)$ 出发解得 $\Delta u=\Delta v=1/2$，交点为 $(1/2,1/2)$。本题方程本来就是线性的，所以一步即达解。}

Newton法的初值过远可能发散；近切触时两切矢接近平行，$J$ 可能病态；重叠曲线则具有一族解，不能简单作为孤立交点求解。一次收敛也只得到一个交点，不能据此断言已找到全部交点。

较稳妥的流程是：先用控制点凸包或\textbf{轴对齐包围盒}剔除不可能相交的曲线段；对包围盒重叠的段用de Casteljau细分；重复剔除与细分，直至局部控制折线足够接近曲线；以折线交点估计 $u,v$，再用Newton精化、验证并去重。包围盒不相交可排除交点，包围盒相交只是\textbf{候选}，不是相交证明。所有距离判断都应使用与模型尺度一致的容差。

\section{B样条拟合：从数据点反求控制点}
\sourcepage{第4章}{57--87}

\subsection{插值与逼近的区别}
给定有序数据点 $\bs Q_0,\ldots,\bs Q_m$，需要确定次数、数据参数 $\bar u_j$、节点及控制点。\textbf{插值}要求 $\bs C(\bar u_j)=\bs Q_j$；\textbf{逼近}只要求误差尽量小。精确设计点适合插值；含噪声的密集扫描点若全部插值，可能产生振荡并需要过多控制点，此时常用较少控制点作最小二乘逼近。

\subsection{给数据点分配参数}
设 $d_i=\|\bs Q_i-\bs Q_{i-1}\|$，$i=1,\ldots,m$。常用一族参数化为
\[
 \bar u_0=0,\qquad
 \bar u_j=\frac{\sum_{i=1}^jd_i^e}{\sum_{i=1}^md_i^e},\qquad
 \bar u_m=1.
\]
其中均匀参数化取 $e=0$，直接令 $\bar u_j=j/m$；弦长法取 $e=1$，用折线长近似弧长；向心法取 $e=1/2$，降低长弦和短弦差异的影响。连续重复点应先合并或专门处理，避免相同参数造成重复方程。

\ex{讲义数据的参数化}{给定 $(0,0),(1,2),(3,4),(4,0)$，弦长分别为 $\sqrt5,\sqrt8,\sqrt{17}$，总长约为 $9.1876$，因此
\[
 \begin{array}{c|cccc}
 \text{方法}&\bar u_0&\bar u_1&\bar u_2&\bar u_3\\\hline
 \text{均匀}&0&1/3&2/3&1\\
 \text{弦长}&0&0.2434&0.5512&1\\
 \text{向心}&0&0.2871&0.6101&1
 \end{array}
\]
弦长的平方根才是向心法累加的量。先算相邻距离，再归一化，避免把坐标差平方和直接当作弦长。}

\reviewfig[.81]{ch04-parameterization.pdf}{同一组数据点采用不同参数化，插值曲线形状也会改变。第4章讲义，PDF第68页。}{fig:parameterization}

参数化改变的是点与参数的对应关系；只有重新求得控制点后，才能看到相应的拟合曲线。均匀分配在数据间距差异大时容易出现凸起或回环，弦长和向心通常更稳妥，但都不是对任意数据的无振荡保证。

\subsection{插值方程与节点平均法}
最基本的全局插值令控制点数等于数据点数，即 $n=m$。固定参数、次数和节点后，有
\[
 A_{ji}=N_{i,p}(\bar u_j),\qquad A\bs P=\bs Q.
\]
$x,y,z$ 三个坐标共用同一个系数矩阵，分别作为右端项求解。端点夹持时，直接得到 $\bs P_0=\bs Q_0,\bs P_m=\bs Q_m$。

选取端点各重复 $p+1$ 次，内部节点用数据参数的连续平均值：
\[
 t_0=\cdots=t_p=0,\quad t_{m+1}=\cdots=t_{m+p+1}=1,
 \qquad t_{j+p}=\frac1p\sum_{i=j}^{j+p-1}\bar u_i,
\]
其中 $j=1,\ldots,m-p$。对严格递增且满足常规插值配置的数据参数，该选择有助于保证插值矩阵可逆。\textbf{不能脱离参数分布任意均匀放置内部节点}，否则某些基函数可能在全部数据参数处都为0，矩阵就会缺列秩。

\ex{三次插值的节点数核对}{6个数据点有 $m=5$，取 $p=3$，共需 $m+p+2=10$ 个节点。若参数为
\[
 0,\quad1/4,\quad1/3,\quad2/3,\quad3/4,\quad1,
\]
则两个内部节点为 $t_4=(1/4+1/3+2/3)/3=5/12$、$t_5=(1/3+2/3+3/4)/3=7/12$，故
\[
 U=(0,0,0,0,5/12,7/12,1,1,1,1).
\]
端点重数4对应三次，而不是四次。}

若取 $n>m$，仅凭插值条件通常不足以唯一确定全部控制点，还需端部导数、光顺准则等附加条件。公式 $\bs P=A^{-1}\bs Q$ 只表达可逆时的数学关系；实际计算应解线性方程组，无须显式求逆，也不应把一般插值矩阵当作对称正定矩阵。

\subsection{最小二乘逼近}
当 $n<m$ 时，给定参数下的拟合通常为超定问题。若强制首末端点吻合，取 $\bs P_0=\bs Q_0,\bs P_n=\bs Q_m$，定义
\[
 \bs R_j=\bs Q_j-N_{0,p}(\bar u_j)\bs P_0
                    -N_{n,p}(\bar u_j)\bs P_n,\quad j=1,\ldots,m-1.
\]
令 $A_{ji}=N_{i,p}(\bar u_j)$，$i=1,\ldots,n-1$，最小化
\[
 E=\sum_{j=1}^{m-1}\left\|\bs R_j-\sum_{i=1}^{n-1}A_{ji}\bs P_i\right\|^2
   =\|\bs R-A\bs P\|_F^2.
\]
对控制点坐标求导，得到正规方程
\[
 A^{\mathsf T}A\bs P=A^{\mathsf T}\bs R.
\]
只有 $A$ 满列秩时，$A^{\mathsf T}A$ 才正定并给出唯一解。实际数值计算可用QR或SVD，减少病态正规方程带来的误差。

内部节点数为 $n-p$，比插值情形少；应结合数据参数分布设置。讲义给出的一种平均分配法是令
\[
 d=\frac{m+1}{n-p+1},\quad i=\lfloor jd\rfloor,\quad
 \alpha=jd-i,\quad t_{p+j}=(1-\alpha)\bar u_{i-1}+\alpha\bar u_i,
\]
其中 $j=1,\ldots,n-p$。它在数据参数之间插值分配节点，仍须检查秩和数据的有效性，不能把节点选择本身当作无条件的数值稳定性保证。

\ex{一个未知控制点的逼近}{对二次Bézier取端点 $\bs P_0=(0,0),\bs P_2=(2,0)$。希望在参数 $1/4,1/2,3/4$ 逼近数据 $(1/2,1),(1,1),(3/2,1)$。由对称性，$P_{1x}=1$；对 $P_{1y}=h$，三个预测高度为 $3h/8,h/2,3h/8$，故
\[
 E(h)=2(3h/8-1)^2+(h/2-1)^2.
\]
令 $E'(h)=0$，得 $h=40/17$。这使\textbf{总平方误差}最小，但三个预测高度仍不都等于1。}

\subsection{参数校正与几何距离}
固定 $\bar u_j$ 后最小化的是 $\|\bs Q_j-\bs C(\bar u_j)\|$，不一定是点到曲线的最短距离。若要逼近几何距离，可以交替进行两步：固定参数求控制点；固定曲线，将各参数移动到数据点的最近点附近。

对一个数据点，令 $f(u)=\tfrac12\|\bs C(u)-\bs Q\|^2$。内部最近点的必要条件为
\[
 g(u)=(\bs C(u)-\bs Q)\cdot\bs C'(u)=0,
\]
即连接矢量垂直于切线。Newton校正为
\[
 \Delta u=-\frac{(\bs C-\bs Q)\cdot\bs C'}
 {\|\bs C'\|^2+(\bs C-\bs Q)\cdot\bs C''},\qquad
 u_{\mathrm{new}}=u+\Delta u.
\]
还应比较区间端点、限制参数范围，并防止多个局部极值和零分母。\textbf{只校正参数不改变曲线几何形状}；重新固定校正后的参数并求控制点，才会更新拟合曲线。迭代可按误差改善幅度或最大迭代数停止。

\Needspace{13cm}
\section{复习速查与自测}
\begin{itemize}
 \item Hermite给端点和端部导矢；Bézier给控制点；B样条还需次数与节点；NURBS再增加权重。
 \item Bézier $n+1$ 个控制点对应最高 $n$ 次；B样条的控制点数与次数可独立选择。
 \item 三次B样条 $p=3,k=4$，单内部节点通常为 $C^2$；二重节点通常为 $C^1$，三重节点通常为 $C^0$。
 \item 升阶、正确插入节点、de Casteljau细分均保持原曲线；移动控制点一般改变曲线。
 \item 凸包或包围盒分离可剔除交点；重叠不能证明相交。Newton局部收敛不保证找到所有交点。
 \item 数据点不等于控制点；插值经过数据点不表示经过所有控制点。最小二乘残差也不自动等于几何最短距离。
\end{itemize}

先独立作答，再核对下列提示。
\begin{enumerate}
 \item 三次Bézier的 $\bs P_0=(0,0),\bs P_1=(1,2)$，起点导矢是多少？\textit{答：$(3,6)$。}
 \item 两段三次曲线接点两侧控制边共线就能保证 $C^1$ 吗？\textit{答：不能；同向平行只对应非退化情形的 $G^1$，$C^1$ 还要求同一参数尺度下导矢相等。}
 \item 二次曲线的一个内部二重节点通常保证几阶连续？\textit{答：$C^0$。}
 \item 8个控制点定义三次B样条，需要多少个节点？\textit{答：$8+3+1=12$ 个。}
 \item 将全部NURBS权重由 $w_i$ 改成 $3w_i$ 会怎样？\textit{答：曲线不变，分子分母同乘3。}
 \item 为什么不能对有噪声的扫描点一律采用精确插值？\textit{答：可能将噪声变成曲线振荡，控制点数量也会过大。}
 \item 参数校正结束而控制点尚未更新，曲线形状改变了吗？\textit{答：没有；只改变了数据点与曲线上点的对应。}
\end{enumerate}
